\documentclass[DIV10, parskip=full*, headings=big, twopage, 10.5pt, cleardoublepage=plain,
   			bibliography=totoc, listof=totoc, BCOR=12mm,
   			numbers=noenddot]{scrbook}

\usepackage[ngerman]{babel}

\usepackage[T1]{fontenc} % for correct guillemets, http://en.wikibooks.org/wiki/LaTeX/Internationalization#German
\usepackage{lmodern} % to avoid rastered fonts, http://www.komascript.de/node/699

\usepackage[babel,german=quotes]{csquotes}
\usepackage[style=alphabetic,backref=true,minnames=3,abbreviate=false]{biblatex}

\usepackage{booktabs}

\usepackage{typearea}

\usepackage{nomencl}

\usepackage{xspace}

\makenomenclature

%Schriften
\usepackage[sc]{mathpazo}
\linespread{1.05}
\usepackage{microtype}
\usepackage{amssymb}

\newenvironment{quoteit}
{\begin{quote}\itshape}
{\end{quote}}

\usepackage{changebar}

% Palatino als Brotschrift
\fontfamily{ppl}\fontseries{l}\selectfont

% AvantGarde für Überschriften
\renewcommand{\sfdefault}{pag} % default für alles serifenlose --> AvantGarde
\setkomafont{subtitle}{\fontfamily{pag}\fontseries{m}\selectfont}
\setkomafont{subject}{\fontfamily{pag}\fontseries{m}\selectfont}
\setkomafont{pagehead}{\sffamily\footnotesize}
\setkomafont{pagenumber}{\sffamily\footnotesize}

\usepackage{setspace}

\usepackage{xcolor}
\definecolor{darkgray}{gray}{0.2} 
\color{darkgray}
\definecolor{lightgrey}{rgb}{0.9, 0.9, 0.9} 
\definecolor{txtgreen}{rgb}{0.549, 0.749, 0.149}

\usepackage{amssymb, amsmath}
\usepackage[final]{pdfpages}

\usepackage{subfigure}
\usepackage{graphicx}
\usepackage{wrapfig}

\usepackage{flafter}
\usepackage{multicol}
\usepackage{multirow}
\usepackage{tabulary}
\usepackage{listings} % zum Einbringen von Code mit Syntaxhightlighting
\usepackage{blindtext}
\usepackage[perpage, ragged]{footmisc}

\definecolor{darkred}{rgb}{0.85, 0.3, 0.4}
\usepackage{courier}
\usepackage{etoolbox}
\lstset{
		basicstyle=\small\ttfamily,
		keywordstyle=\color{blue},
		stringstyle=\color{darkred},
		commentstyle=\tiny\color{gray}
}
% voreingestellte Programmiersprache: Java
\lstset{
        language=Java, 
        breaklines=true, 
        breakautoindent=true, 
        breakatwhitespace=true, 
        numbers=left, 
        numberstyle=\tiny,
        extendedchars=true,
        rulesepcolor=\color{red},
        frame=lines,
        keywordstyle=\color{blue},
        commentstyle=\color{gray},
        stringstyle=\color{darkred},
        emph={double,bool,int,unsigned,char,true,false,void},
				emphstyle=\color{blue},
				emph={Assert,Test},
				emphstyle=\color{red},
				emph={[2]\using,\#define,\#ifdef,\#endif}, emphstyle={[2]\color{blue}},
		basicstyle=\small\ttfamily,
		showstringspaces=false,
		commentstyle=\tiny}
		
\usepackage[
   	hyperfootnotes=false,
   	bookmarks,
   	pdfpagelabels=true,
   	plainpages=false,
   	breaklinks,
   	unicode,
   	pdfencoding=auto]{hyperref}
	
\definecolor{lc}{cmyk}{0.6,0.16,1,0.015} % for screen output
\newcommand{\printoutput}{\definecolor{lc}{cmyk}{0,0,0,1}} % for print output

\hypersetup{
	pdftoolbar=true,
	bookmarksopen,
	bookmarksnumbered=false,
	bookmarksopenlevel=1,
	pdfdisplaydoctitle,	
	colorlinks=true,
	linkcolor=lc,
	citecolor=lc,
	filecolor=lc,
	menucolor=lc,
	urlcolor=lc,
	pdfauthor={$author$},
	pdftitle={$title$},
}

\usepackage[automark, footsepline]{scrpage2}

%%%%%%%%%% Tabellenanpassungen %%%%%%%%%%%%%%%%%%
\setlength{\tabcolsep}{3pt}
\renewcommand{\arraystretch}{2}

\usepackage{colortbl}
\usepackage{ltxtable, array}
\usepackage{longtable}
\usepackage{caption}
\captionsetup{font=small}
\captionsetup{labelfont=bf}
\renewcommand{\tabularxcolumn}[1]{m{#1}} % Zellen von Tabellen immer vertikal zentrieren

%%%%%%%%%%%%%% Listen %%%%%%%%%%%%%%%%%%%%%
\usepackage{mdwlist}
\usepackage{setspace}
\renewcommand{\labelitemi}{\textperiodcentered}%{\sqbullet} %{\textperiodcentered}
\usepackage{enumitem}
\setlist{noitemsep}
\renewcommand{\labelenumi}{(\arabic{enumi})}

%%%%% Aufzählungsstyle %%%%%%%%
\newcommand{\myitem}[1]{\textsf{\textbf{#1}}.\hspace{1em}}

%%%%%%%%%% Pagina und Kolumnentitel %%%%%%%%%%%%%
\pagestyle{scrheadings}
\rohead{\headmark}
\ofoot{\pagemark}
\setheadsepline{0.5pt}
\setfootsepline{0pt}

\setcounter{secnumdepth}{3} % Tiefe der Überschriften, die Nummern bekommen
\setcounter{tocdepth}{2}    % Inhaltsverzeichnis bis Tiefe 3

%%%%%%%%%%%%%%%%%%%% TODO Notes %%%%%%%%%%%%%%%%%%%%%%%%
\usepackage[colorinlistoftodos]{todonotes} % TODO notes

\usepackage{fancyvrb}
\usepackage{fontspec}

\newcommand{\citet}{\cite}
\newcommand{\citep}{\cite}

% \usepackage[inline]{showlabels}
\printoutput

% Bib-File
\bibliography{bibliography.bib}
\begin{document}


\begin{titlepage}
\begin{figure}[htbp]
  \begin{minipage}[t][3cm][c]{\textwidth/2+1.7cm}
	
		\includegraphics[height=3.5em]{figures/tud.pdf}  

  \end{minipage}
  \begin{minipage}[t][3cm][c]{\textwidth/2-1.7cm}
	\begin{flushright}  	
  	\includegraphics[height=3.5em]{figures/logo.pdf}
	\end{flushright}
  \end{minipage}
\end{figure}


\vspace{10em}
\large{\textsf{$type$}}
\newline
\Huge{\textbf{\textsf{$title$}}}

\vspace{\fill}

\normalsize 

\begin{flushleft}

\begin{figure}[b]
\centering
\begin{minipage}[b]{\textwidth/2-0.1cm+0.85cm}
\begin{flushleft}
$date$\\ \vspace{1em}
$author$\\ \vspace{0.25em}
Matr.-Nr.: $matrikel$
\end{flushleft}
\end{minipage}
\begin{minipage}[b]{\textwidth/2-0.1cm-0.85cm}
\begin{flushleft}
\textbf{\textsf{Betreuer}}\\ \vspace{0.25em}
$advisor$\\
\vspace{1em}
\textbf{\textsf{Verantwortl. Hochschullehrer}}\\ \vspace{0.25em}
$teacher$
\end{flushleft}
\end{minipage}
\end{figure}

\end{flushleft}

\end{titlepage}

\pagenumbering{roman}

\thispagestyle{empty}
\vspace*{\fill}

\section*{Erklärung}
Hiermit erkläre ich, dass ich die vorliegende Arbeit mit dem Titel

\vspace*{1em}
\begin{large}
\textsf{\textbf{$title$}}
\end{large}
\vspace*{1em}

unter Angabe aller Zitate und nur unter Verwendung der angegebenen Literatur
und Hilfsmittel selbstständig angefertigt habe.

Dresden, den $date$\\
\vspace*{3em}

\emph{$author$}

\linespread{1.13}

\pdfbookmark[1]{Inhaltsverzeichnis}{toc}
\tableofcontents					% create a toc
%\addcontentsline{toc}{chapter}{Inhaltsverzeichnis}
% wenn eine Seite beim Inhaltsverzeichnis fehlt (also zB drei statt vier Seiten lang ist):
%\thispagestyle{empty}
\cleardoublepage

\linespread{1}
%\pagenumbering{arabic}

% bei vierseitigen Inhaltsverzeichnis erst im Kapitel 1
\pagenumbering{arabic}

%%%%%%%%%%%%%%%%%%%% Struktur %%%%%%%%%%%%%%%%%%%%%%%%

$body$

\thispagestyle{empty}
%\nocite{*}

%%%%%%%%%%%%%%%%%%%%%%%%%%% ANHANG %%%%%%%%%%%%%%%%%%%%%%%%%%%%%%%%%%%
\begin{appendix}
\pagenumbering{roman}
\clearpage
\setcounter{page}{1}

$for(include-after)$
$include-after$
$endfor$

%%%%% ABBILDUNGEN %%%%%
\renewcommand{\listfigurename}{Abbildungen}
\listoffigures

%%%%% CODE %%%%%
\renewcommand\lstlistlistingname{Programmcode}
\lstlistoflistings

%%%%% TABELLEN %%%%%
\renewcommand\listtablename{Tabellen}
\listoftables

%%%%% BIBLIOGRAFIE %%%%%%
\printbibliography[title=Literatur]

\end{appendix}
\end{document}
